% OWNER: sections/s1_intro.tex. Edit only this file.
\seclabel
\claim{\secnum{1}{sec:intro}\quad Without a published API, we measured the server through its web client's requests}

\lead{The competition takes orders only through a web page, so a program that trades needs a wrapper in place of the browser.} Road to Fund Manager is a paper-trading competition run by Timefolio Asset Management; participants enter tickers and target weights on a web page. An algorithmic strategy has to log in, read balances, follow positions as they change and place orders on its own. The wrapper handles that communication so that the strategy code never deals with tokens, sockets or date formats. The wrapper sends the HTTP requests the page itself sends, and Section~\ref{sec:design} compares this with driving the page in a browser.

\lead{The organiser publishes no API, so every assumption the wrapper makes has to come from observation.} That covers how authentication works, how long tokens last, how dates are read and what the live messages contain. Guessing is risky because the server rarely complains: as Section~\ref{sec:queries} shows, a wrong date returns another day's data with HTTP 200, so a wrong assumption produces wrong numbers without an error.

\lead{The organiser allows automated orders within 2 limits.} In reply to our enquiry on 24 Nov 2026 [5], the organiser stated that it provides no official API and does not prohibit automating orders with scripts or browser-automation tools (web drivers). An IP address that places 100 or more orders in a day may be blocked, and so may bursts of traffic that put a heavy load on the system; the organiser recommends a delay between requests but gave no figure for it. We kept a copy of the reply. We ran no load or rate-limit tests. Our heaviest traffic came from the timing study of Section~\ref{sec:transport}: runs of 5 sequential requests, about 170\ms{} per run, repeated 10 times at 1\,s intervals. The load threshold is unknown.

\lead{The page's JavaScript and its recorded traffic gave every request path and feed method the wrapper uses.} After logging in we extracted API paths and live-feed method names from the JavaScript the site serves, and recorded WebSocket messages in the browser's developer tools. We then reproduced each request with curl to read headers, cookies and response bodies; for the WebSocket, which curl cannot handle, we implemented the SignalR protocol in Python without a client library.

\lead{Where a behaviour had more than one explanation, requests designed to separate the explanations settled it.} We wrote down the competing hypotheses and chose requests whose replies differ under each. For the condition that triggers a delta, 11 requests on 2 connections left 1 of 3 candidate rules standing (Section~\ref{sec:feed}); for the server's share rounding, 1 order worth 1.76 shares separated rounding down from rounding to the nearest share (Section~\ref{sec:oneshare}).

\lead{The report runs from measurement to design to live tests.} Section~\ref{sec:layers} measures 5 protocol layers: transport (\ref{sec:transport}), authentication (\ref{sec:auth}), queries (\ref{sec:queries}), the live feed (\ref{sec:feed}) and orders (\ref{sec:orders}). The first 4 were measured mostly before the 09:00 open and rechecked after it where the trading session could change the result; orders were measured during the session. Section~\ref{sec:design} turns the measurements into the wrapper's goals, structure and operating rules. Section~\ref{sec:live} tests the finished wrapper against the live server: start-up, feed and orders (\ref{sec:wrappertests}), the check that keeps orders from rounding to 0 shares (\ref{sec:oneshare}), the competition's investment guidelines (\ref{sec:guidelines}), and the limits of these tests with the questions still open (\ref{sec:validity}). \figref{setup} gives the test setup; Appendix~\ref{app:glossary} defines the terms and Appendix~\ref{app:decisions} lists the design decisions with their evidence.

% OWNER: figs/setup.tex. Edit only this file.
\fig{setup}{All tests used one account and one client on 24 Nov 2026}[Package, test window, client, account, server and load behind every measurement in this report]{%
{\tabfont\begin{tabularx}{\linewidth}{>{\RaggedRight\arraybackslash}p{16mm} >{\RaggedRight\arraybackslash}X}
\thd{Item} & \thd{Value}\\
Package & \code{pytimefolio} v0.1.0, MIT licence [8]; \href{https://github.com/chlanstjr0301/pytimefolio}{\code{github.com/chlanstjr0301/pytimefolio}}\\\rowsep
Time & 05:38 to 10:34 KST on 24 Nov 2026; the market opened at 09:00. The server's business day (\code{currBizDay}) changed from 23 Nov to 24 Nov in this window\\\rowsep
Client & WSL (Ubuntu) on Windows 11; curl 8 with HTTP/2; Python 3.14, httpx 0.28.1 with h2, websockets 17.4. PC clock 3\,s ahead of the server's \code{Date} header\\\rowsep
Account & One participant account; one live portfolio (RFM round 13, 12 Oct to 30 Nov 2026)\\\rowsep
Server & \code{https://contest.timefolio.net} (221.168.38.40); TLS 1.3 with a Let's Encrypt certificate; REST under \code{/api}, live feed at \code{/hubs/main}\\\rowsep
Load & Timing runs as a burst: 5 back-to-back requests in about 170\ms. Live-feed requests every 3\,s. Orders in the day: 9 new orders and 7 stops, against the organiser's threshold of 100 orders\\
\end{tabularx}}}{Network timings hold for this client and time of day only.}{DART Quant Team}

